\documentclass[11pt]{article}
\usepackage[margin=0.9in]{geometry}
\usepackage{amsmath,amssymb,amsthm,booktabs,graphicx,microtype,needspace}
\usepackage[hidelinks]{hyperref}
\newtheorem{theorem}{Theorem}
\newtheorem{lemma}{Lemma}
\newcommand{\Sph}{\mathbb S^d}
\newcommand{\dist}{\operatorname{dist}_{\Sph}}
\setlength{\parindent}{0pt}
\setlength{\parskip}{0.45em}

\title{One Zone Disproves the Printed Zone--Inradius Conjecture}
\author{Run \texttt{fa\_banach\_001}, lane 2}
\date{23 July 2026}

\begin{document}
\maketitle

\textbf{Status.} Candidate counterexample, likely valid, for human review.

\textbf{Outcome.} Conjecture~8 in Polyanskii's \emph{A cap covering
theorem} is false as literally printed.  In every dimension, any single
nondegenerate zone is already a counterexample: the conjectured left-hand
side equals $\pi$, whereas the claimed lower bound is $2\pi$.

\section{Source question}

Let $\Sph\subset\mathbb R^{d+1}$ be the unit sphere.  The source defines a
zone of width $2\alpha$ as the set of points at spherical distance at most
$\alpha$ from a great sphere \cite[p.~2]{Polyanskii}.  Its final open problem
is the following.

\begin{center}
\includegraphics[width=\textwidth]{figures/open_problem_crop.png}
\end{center}

\begin{quote}
\textbf{Conjecture 8.}
Let \(Z_1,\ldots,Z_n\subset S\) be zones of width
\(\beta_1,\ldots,\beta_n\), respectively. If
\(S\setminus\bigcup_{i=1}^n Z_i\) consists of convex connected components
\(L_1,\ldots,L_{2m}\) with inradius
\(\gamma_1,\ldots,\gamma_{2m}\), respectively, then
\[
 \beta_1+\cdots+\beta_n+\gamma_1+\cdots+\gamma_{2m}\ge 2\pi.
\]
\end{quote}

The paper attributes the conjecture to ``Maxim Didid''; its source label is
\texttt{conjecture:didin}.  We refer to it simply as the printed
zone--inradius conjecture.

\Needspace{13\baselineskip}
\section{Claimed contribution}

\begin{theorem}\label{thm:counterexample}
Fix $d\ge 1$ and $0<\beta<\pi$.  Let $Z\subset\Sph$ be any zone of width
$\beta$.  Then $\Sph\setminus Z$ has exactly two convex connected components,
each of spherical inradius
\[
 \gamma=\frac{\pi-\beta}{2}.
\]
Consequently, with $n=1$ and $m=1$, the left-hand side of Conjecture~8 is
\[
 \beta+\gamma+\gamma=\pi<2\pi.
\]
Thus Conjecture~8 is false as stated.
\end{theorem}

\section{Proof intuition}

A zone is a symmetric belt around an equator.  If its half-width is
$\beta/2$, then each of the two polar caps left over has angular radius
$\pi/2-\beta/2$.  The belt width plus the radii of both caps therefore
telescopes:
\[
 \beta+2\left(\frac{\pi}{2}-\frac{\beta}{2}\right)=\pi.
\]
This is independent of both the width and the dimension, and is only half
of the conjectured lower bound.

\section{Proof}

Choose a unit vector $u\in\mathbb R^{d+1}$ normal to the great sphere around
which $Z$ is centered.  Put $a=\beta/2$.  For $x\in\Sph$, the spherical
distance from $x$ to the great sphere $u^\perp\cap\Sph$ is
\[
 \arcsin |\langle x,u\rangle|.
\]
Since the source's zones are closed, it follows that
\[
 Z=\{x\in\Sph:|\langle x,u\rangle|\le \sin a\}.
\]
Hence
\[
 \Sph\setminus Z=L_+\mathbin{\dot\cup}L_-,
 \qquad
 L_\pm=\{x\in\Sph:\pm\langle x,u\rangle>\sin a\}.
\]
Because
\[
 \sin a=\cos\left(\frac{\pi}{2}-a\right),
\]
and $\dist(x,u)=\arccos\langle x,u\rangle$, we have
\[
 L_+=\left\{x\in\Sph:\dist(x,u)<\frac{\pi}{2}-a\right\},
\]
with the analogous formula for $L_-$ centered at $-u$.  Thus $L_+$ and
$L_-$ are antipodal open spherical caps of radius
\[
 r=\frac{\pi}{2}-a=\frac{\pi-\beta}{2}.
\]
Since $0<r<\pi/2$, each cap is geodesically convex and connected, and there
are no other components.

For completeness, the inradius of either open cap is exactly $r$.  Every
concentric cap of radius $r'<r$ is contained in it, so its inradius is at
least $r$.  Conversely, if a spherical cap of radius $\rho$ is contained in
$L_+$ and has center $y$, the triangle inequality along a minimizing
geodesic from $u$ through $y$ gives
\[
 \rho\le r-\dist(u,y)\le r.
\]
Equivalently, if inradius is defined as a supremum for open domains, the
supremum is $r$.  Therefore $\gamma_1=\gamma_2=r$, and the quantity in the
conjecture is
\[
 \beta+\gamma_1+\gamma_2
 =\beta+2\left(\frac{\pi-\beta}{2}\right)
 =\pi<2\pi.
\]
All of the hypotheses of Conjecture~8 hold with $n=1$ and $2m=2$, completing
the proof.

\section{Hypothesis audit}

\begin{center}
\begin{tabular}{@{}lll@{}}
\toprule
Printed requirement & One-zone construction & Check \\
\midrule
Zones $Z_1,\ldots,Z_n$ & one zone $Z$, $n=1$ & yes \\
Widths $\beta_i$ & $\beta_1=\beta\in(0,\pi)$ & yes \\
Even number of components & $L_+,L_-$, so $2m=2$ & yes \\
Components connected & open caps & yes \\
Components convex & radii $<\pi/2$ & yes \\
Component inradii & $\gamma_1=\gamma_2=(\pi-\beta)/2$ & yes \\
Claimed inequality & $\pi\ge 2\pi$ & false \\
\bottomrule
\end{tabular}
\end{center}

\section{Robustness, interpretation, and limitations}

The counterexample does not depend on a limiting or degenerate zone: every
$\beta\in(0,\pi)$ works.  For example, $\beta=\pi/2$ leaves two caps of
inradius $\pi/4$, and the conjectured sum is again $\pi$.

The open/closed convention does not change the calculation.  The source
defines the zone as the closure of the intersection with an open plank, so
its complementary components are open.  Their inradius, in the standard
supremal sense for domains, is $(\pi-\beta)/2$.  Passing to closures yields
closed caps with the same numerical inradius.

The factor-of-two discrepancy strongly suggests that the displayed
constant, one of the summands, or a normalization was misstated.  This
packet disproves only the literal printed statement; it does not assert
which correction was intended, and it does not address Conjecture~7 in the
same paper.

\clearpage
\section{Verification and novelty search}

The script \texttt{code/verify\_one\_zone.py} checks the exact symbolic
identity and several sample widths.  It runs with:
\begin{verbatim}
conda run --no-capture-output -n sandbox python \
  code/verify_one_zone.py
\end{verbatim}

Before promotion, the run's cheap title, arXiv-id, and keyword indexes were
searched.  Targeted web searches for the exact displayed inequality, the
paper title together with ``correction'' and ``erratum,'' and the
zone/inradius keywords found no later correction or resolution as of
23 July 2026.  The local arXiv source and PDF both contain the $2\pi$
constant.  This is a bounded novelty check, not a comprehensive literature
review.

\section{Human review recommendation}

A reviewer in discrete or spherical geometry should first confirm the
source's intended convention for ``inradius'' of an open component and ask
the author whether $2\pi$ is a typographical or normalization error.  Under
the definitions printed in the paper, the proof above is complete.

\begin{thebibliography}{1}
\bibitem{Polyanskii}
A. Polyanskii,
\emph{A cap covering theorem},
Combinatorica \textbf{41} (2021), 695--702;
\href{https://arxiv.org/abs/2006.02192}{arXiv:2006.02192}.
\end{thebibliography}

\end{document}
